\section{The simultaneous $r,j$ opening: a harmonic full-event frame}
\label{sec:r13-c1-simultaneous-rj}

\begin{theorem}
Use the original physical native C1 source, with $d=1$ and $x=k^2>0$
infinitesimal. Suppose $\eta,\zeta$ are positive bounded units and
$r$ and the actual $j=\eta(1+xh)-1+r$ are both positive
infinitesimals. Then the complete original parent is tensor SOS for
every independent projective test and every state. Literal $\beta=0$
is retained. No relative rate among $x,r,j$ or between source and
testing parameters is imposed.
\end{theorem}

All scalar and Gram identities in this section use the normalized
parent $F/4$; multiplication of every square root by $2$ restores $F$.

The essential change is to combine the whole positive $A,D$ events
before choosing the scalar graph. The weighted $uvX$ term is retained
inside a hard root. Solving for $p=uW$, rather than for $q=vZ$,
then gives a polynomial inverse with a pure-source unit denominator.

\paragraph{Source bounds and the harmonic frame.}
Write $U=u^2$, $V=v^2$, $D_u=1+U$, $D_v=1+V$,
$\theta=1-h=x\pi$, $B_x=\beta+1/2$, and
\[
 T=1+xh,\qquad e=\eta T=1-r+j,\qquad y=x\eta\zeta .
 \tag{HO1}
\]
The bounded-$\eta$ source bound gives $\pi=O(1)$.
Here $e,\eta,h,1-r\sim1$. Exact native numerator cancellation gives
\[
 \begin{gathered}
 C-C_* =O(x),\qquad Ar-e/2=O(x),\qquad Dj-1/2=O(x),\\
 C_*={1+j\over2(1-r)},\qquad A_*={e\over2r},\qquad D_*={1\over2j},\\
 m/x\sim1/(2\zeta),\qquad
 \alpha/x^2\sim\zeta(1+\zeta)/2,\\
 B_x\theta/x=O(1),\qquad
 q_U/x^2=O(1),\quad q_U=\alpha+(m+B_x)\theta^2 .
 \end{gathered}
 \tag{HO2}
\]
The checker cancels the exact rational formulas before setting
$x=r=j=0$: every denominator in the seven displayed bounded
coquotients factors there into constants and powers of the retained
unit $\zeta$. This proves boundedness at arbitrary simultaneous
source rates, rather than just nonvanishing at a selected boundary.
In particular $A,D>0$ and both are unbounded.
Also $\theta/\sqrt\alpha$ and $B_x\theta^2/\alpha$ are bounded,
including $\beta=0$.

Define pure-source positive quantities
\[
 \mathsf H=A+D,\qquad \rho={D\over\mathsf H},\qquad
 \mathsf K={AD\over\mathsf H},\qquad
 \delta={Ar+Dj\over\mathsf H}.
 \tag{HO3}
\]
Both $\mathsf H$ and $\mathsf K$ are unbounded,
$0<\rho<1$, and
$\mathsf K^{-1}=A^{-1}+D^{-1}=O(r+j)$.
Put $p=uW$, $q=vZ$, $J_0=p+uvX-q$, and use the full-event roots
\[
 \begin{aligned}
 E&=\sqrt{\mathsf K}\{J_0-(r-j)p\},\\
 N&=\sqrt{\mathsf H}\{Y-\rho J_0-\delta p\},\\
 N_x&=\sqrt{B_x}(X+huZ),\qquad
 \ell=E/\sqrt{\mathsf K},\quad n=N/\sqrt{\mathsf H},
 \quad \xi=N_x/\sqrt{B_x}.
 \end{aligned}
 \tag{HO4}
\]
All three vanish at the actual source contact. Their exact inverse is
\[
 \begin{aligned}
 p&={(1+hU)q-uv\xi+\ell\over e},&
 Y&=rp+\rho\ell+n,&X&=-huZ+\xi,\\
 J_t&=-jp+(1-\rho)\ell-n,&J_t&=p+uvX-Y-q .
 \end{aligned}
 \tag{HO5}
\]
The denominator $e$ is pure source. In particular no inverse of
$1+hU$ occurs, even if the testing and source rates are unrelated.
The useful exact weighted identities are
\[
 {A\rho^2\over\mathsf K}=\rho,\qquad
 {A\over\mathsf H}=1-\rho,\qquad
 {D(1-\rho)^2\over\mathsf K}=1-\rho,\qquad
 {D\over\mathsf H}=\rho .
 \tag{HO6}
\]

\paragraph{The complete hard block and its reference cancellation.}
The $A,D$ base has the exact identity
\[
 \begin{split}
 AY^2+DJ_t^2={}&N^2+E^2
 +2p\left\{{Ar+Dj\over\sqrt{\mathsf H}}N
                 +\sqrt{\mathsf K}(r-j)E\right\}\\
 &+(Ar^2+Dj^2)p^2 .
 \end{split}
 \tag{HO7}
\]
Both displayed cross coefficients and the last coefficient are
infinitesimal: the squared $E$ cross coefficient is $O(r+j)$
and the squared $N$ cross coefficient is $O(\mathsf H^{-1})$.
For bounded testing ratio $t=v/u$, (HO5) has bounded coefficients
on the hard roots, and the only potentially large $m$ coefficient is
$mt^2/\mathsf K$, which is infinitesimal.
The higher $A$ packet is controlled by (HO6); for example
\[
 {Ar^2\over\mathsf K}=r^2+{Ar\over Dj}rj=O((r+j)^2),
 \qquad {Ar\rho\over\mathsf K}=r .
 \tag{HO8}
\]
Consequently its hard coefficients carry an additional $U$ or $V$.
The bounded $C$ and half-core packets have infinitesimal coefficients
on $E,N$, and combining the beta packet with its half-core part gives
the third hard pivot. The full hard residue on $(E,N,N_x)$ is
$\operatorname{diag}(1,1,1)$.

A stronger exact base statement is needed for unbounded $t$.
Set $\widetilde q=q-uvX=ep-\ell$, so that
$J_t=p-Y-\widetilde q$. The reference quadratic form
\[
 \mathcal B_*=
 C_*(p-Y)^2+A_*Y^2+D_*(p-Y-\widetilde q)^2
                    +\tfrac12(p-Y)^2-p\widetilde q
 \tag{HO9}
\]
has no soft $p$ coordinate after substituting
$Y=rp+\rho\ell+n$ and $\widetilde q=ep-\ell$:
\[
 \mathcal B_*=
 (A_*+C_*+\tfrac12)(\rho\ell+n)^2
                 +D_*\{(1-\rho)\ell-n\}^2 .
 \tag{HO10}
\]
This is an exact contact cancellation, checked before taking residues.
Since $A_*/A=1+O(x)$ and $D_*/D=1+O(x)$, (HO6) shows that
its bounded hard Gram on $(E,N)$ has residue $\operatorname{diag}(1,1)$.
Replacing the reference weights by the actual $C,A,D$ creates a
quadratic error with coefficients $O(x)$ in $(p,E,N)$.
For the singular weights this follows from
$(A-A_*)Y^2=O(x)\,AY^2$ and
$(D-D_*)J_t^2=O(x)\,DJ_t^2$, using (HO6),(HO7);
neither an $x/r$ nor an $x/j$ estimate is used.

\paragraph{The exact bounded-ratio scalar circuit.}
At $E=N=N_x=0$ and $Z=1$, let
\[
 w(U)={1+hU\over e},\qquad
 p=w(U)v,\quad Y=rw(U)v,\quad X=-hu .
 \tag{HO11}
\]
Since $w(x)=1/\eta$, this graph agrees exactly with the native
source contact. The source contact value and all first state and
angular jets are therefore inherited by the chain rule.

For compact scalar coefficients set
$a=C(1-r)^2$, $b=Ar^2$, $d=(D+1/2)j^2$,
$c=(\alpha+1/2)r^2$, and
\[
 L_A(U)=1+2h+(2h+h^2)U+h^2U^2 .
 \]
The complete scalar is exactly
\[
 \mathcal S=c_F(V-y)^2+f_A(U,V)A_t^2+f_BB_t^2+
                         f_C(U)C_t^2-f_GG^2 ,
 \tag{HO12}
\]
on the separated roots
$A_t=uv$, $B_t=kv^2/(su)$, $C_t=xs$, $G=kv$,
where $s=\sqrt{\eta\zeta}$. The coefficients are
\[
 \begin{aligned}
 c_F&={2m(h+r)+a+b+c\over e^2},\\
 q_V&={a+b+d\over e^2}-{2m\theta\over e}
                       -{(2\alpha+1)r\over e}-\tfrac12,\\
 f_B&={m\eta\zeta\over e^2x},\qquad
 f_C={\alpha-c_Fy^2+q_UU\over\eta\zeta x^2},\qquad
 f_G=-{q_V+2c_Fy\over x},\\
 f_A&={mV\over e^2}
 \{(h+r)^2+2hr+2hr(h+r)U+h^2r^2U^2\}
       -{2m\theta\over e}(h+r+hrU)\\
 &\quad+{a\over e^2}(2h+h^2U)(1+V)
       +{b\over e^2}L_A(U)(1+V)
       +{d\over e^2}(2h+h^2U)+{c\over e^2}L_A(U)V\\
 &\quad-{(2\alpha+1)r\over e}(1+h+hU)
                              +B_x\theta^2-\tfrac12 .
 \end{aligned}
 \tag{HO13}
\]
The only opening weights here are $Ar^2$ and $Dj^2$, both
infinitesimal. Thus $f_A$ is a bounded polynomial of degree at most
two in $U$ and one in $V$, and all nonconstant coefficients of
$f_C$ are bounded by (HO2). Direct source reduction gives
\[
 c_F\sim\tfrac12,\qquad f_A,f_B,f_C\sim\tfrac12 .
 \tag{HO14}
\]
Exact scalar contact gives
$f_G=f_A(x,y)+f_B+f_C(x)\sim3/2$, without separately estimating
$q_V/x$.

Apply the finite three-circuit lemma with $\lambda=1/2$.
The polynomial Taylor arrays are the monic quotients in (JO13),
now applied to
$f_A(u^2,u^2t^2)+f_B+f_C(u^2)-f_G$.
Their $u$ and mixed arrays are bounded; the $t,t$ array and the
coefficient $t$ first differences are infinitesimal, since every
$t$ dependence at $u=k$ contains $k^2$.
The two exact angular contact jets give the two circuit balances.
This supplies the strict three-circuit Gram and the strict
$V-y$ pivot for the entire scalar, including its negative soft term.

\paragraph{All complete first normals.}
The following formulas come from differentiating the entire original
parent after (HO5), with $Z=1$. Put
$a_E=r/e+\rho$, $b_E=(1-r)/e-\rho$,
$c_E=1-\rho-j/e$. The bounded polynomials are
\[
 \begin{aligned}
 P_N/2={}&m\{wV(1+rU)-\theta U\}
 -CD_v(1-r)w+AD_uD_vrw\\
 &+(D+\tfrac12)jw
             +(\alpha+\tfrac12)D_u(rwV-1),\\
 P_E/2={}&{mV\over e^2}(h+r+hrU)-{m\theta\over e}
 +ma_E\{wV(1+rU)-\theta U\}\\
 &+CD_v(1-r)wb_E+AD_uD_vrwa_E
 -(D+\tfrac12)jwc_E
            +(\alpha+\tfrac12)D_u(rwV-1)a_E,\\
 Z_X={}&{m\theta(1+rU)\over e}+B_x\theta\\
 &-{CD_v(1-r)^2w+AD_uD_vr^2w+(D+\tfrac12)j^2w
                 +(\alpha+\tfrac12)D_ur(rwV-1)\over e}.
 \end{aligned}
 \tag{HO15}
\]
Their only opening-weight factors are $Ar,Dj,Ar^2,Dj^2$.
Every coefficient is bounded uniformly in the source rates and
$0<\rho<1$. The three complete first normals are
\[
 \begin{aligned}
 \mathcal L_N&={vP_N\over\sqrt{\mathsf H}},\\
 \mathcal L_E&={2mvt^2/e^2+vP_E\over\sqrt{\mathsf K}},\\
 \mathcal L_x&={1\over\sqrt{B_x}}\left\{
 2u\theta(m+B_x)-2mwut^2(1+rU)
           \left(1+{V(1+rU)\over e}\right)+2uVZ_X\right\}.
 \end{aligned}
 \tag{HO16}
\]
For either $P=P_N,P_E$, use its exact finite monic decomposition
$P=P_{00}+UP_U+VP_V$.
State contact gives
\[
 \begin{aligned}
 P_{N,00}/x&=-P_{N,U}(x,y)-\eta\zeta P_{N,V}(y),\\
 P_{E,00}/x&=-2(m/x)\eta\zeta/e^2
                         -P_{E,U}(x,y)-\eta\zeta P_{E,V}(y).
 \end{aligned}
 \tag{HO17}
\]
Thus these coquotients are bounded without any new opening inverse.
Allocate $vP=(P_{00}/k)G+u(P_U+t^2P_V)A_t$.
The extra $E$ term allocates on $B_t$ with coefficient
$2mst/(e^2k\sqrt{\mathsf K})$.
The three $X$ terms allocate respectively on $C_t,B_t,A_t$,
with coefficients
\[
 {2u\theta(m+B_x)\over xs\sqrt{B_x}},\qquad
 {-2msw(1+rU)\{1+V(1+rU)/e\}\over k\sqrt{B_x}},\qquad
 {2vZ_X\over\sqrt{B_x}} .
 \tag{HO18}
\]
Every displayed allocation coefficient is infinitesimal.
The coefficient sums vanish at contact by the exact state jets;
their finite $u$ first differences are bounded and their $t$
first differences at $u=k$ are infinitesimal.
This follows directly from the bounded polynomial arrays in
(HO15): every $t$ dependence carries $u$ or $k$.
The two unchanged circuit conductors therefore place every whole
normal on the strict circuit roots with infinitesimal coefficients.
Together with the complete hard block, this proves all bounded $t$,
including infinitesimal $t$, without a missing state normal.

\paragraph{The reciprocal ratio: retain both soft directions.}
If $v/u$ is unbounded, work homogeneously and use seven roots
\[
 E,\quad N,\quad N_x,\quad
 P_m=\sqrt m\,vW,\quad Q_\alpha=\sqrt\alpha Z,\quad
 R_t=uvZ,\quad S_t=uvW .
 \tag{HO19}
\]
The base error following (HO10) has zero Gram residue on this list:
use the exact separated identity
$\sqrt x\,p=\sqrt{x/m}(u/v)P_m$.
This absorbs every $O(x)$ coefficient in $(p,E,N)$, since $x/m$
is a source unit and $u/v$ is a pure-testing infinitesimal.

After removing the actual quadratic base (HO9), with actual $C,A,D$,
and the three entire positive packets
$mM_t^2$, $\alpha D_u(vY-Z)^2$,
$B_xD_v(X+uZ)^2$, the full remaining quadratic is exactly
\[
 \begin{split}
 \mathcal R={}&CV(p-Y)^2+A(U+V+UV)Y^2
                      +\tfrac12D_uVY^2-\theta UqY\\
 &+(h-\tfrac12)Uq^2+\tfrac12h^2U^2q^2
       -uv\xi\{Y+(1+hU)q\}+\tfrac12UV\xi^2 .
 \end{split}
 \tag{HO20}
\]
This identity retains the whole $D$ event in the quadratic base,
so no coefficient containing a bare divergent $D$ remains.
It is checked as one complete homogeneous polynomial.

Here is a finite allocation of every term. The exact identity
\[
 up={(1+hU)R_t-u^2v\xi+uE/\sqrt{\mathsf K}\over e}
 \tag{HO21}
\]
has bounded coefficients and uses no mixed inverse.
The $C$ packet has leading root
$\sqrt C(1-r)S_t$; its other coefficients are
$-\sqrt C\,\rho v/\sqrt{\mathsf K}$ and
$-\sqrt C\,v/\sqrt{\mathsf H}$ on $E,N$, both infinitesimal.
All three $A$ higher packets have zero residue. For example,
\[
 \sqrt A\,uY=\sqrt A\,r(up)+\sqrt\rho\,uE+
                               \sqrt{1-\rho}\,uN,\qquad
 \sqrt A\,vY=\sqrt A\,rS_t+\sqrt\rho\,vE+
                               \sqrt{1-\rho}\,vN .
 \tag{HO22}
\]
Here $\sqrt A\,r=O(\sqrt r)$ and (HO8),(HO21)
control every coefficient; the $uvY$ packet is $v$ times the first.
The half-$VY^2$ term is smaller by the same identities.
The $-\theta UqY$ term is $-\theta R_t(uY)$ and has zero residue.
The next two terms are $(h-1/2)R_t^2$ and
$h^2U R_t^2/2$.
Finally, the cross term is
$-\xi v(uY)-\xi v(1+hU)R_t$ and all its coefficients,
as well as those of $UV\xi^2/2$, are infinitesimal.
Thus this allocation exhausts the complete quadratic (HO20);
its only nonzero soft residues are $1/2$ on each of $R_t,S_t$.

The three retained positive packets have leading roots
$P_m,-Q_\alpha,N_x$:
\[
 \begin{aligned}
 \sqrt m\,M_t&=P_m+\sqrt m\,uvY-\sqrt m\,\xi
                  -\sqrt{m/\alpha}\,\theta uQ_\alpha,\\
 \sqrt{\alpha D_u}(vY-Z)
             &=\sqrt{D_u}(\sqrt\alpha\,vY-Q_\alpha),\\
 \sqrt{B_xD_v}(X+uZ)
             &=\sqrt{D_v}\{N_x+
                      \sqrt{B_x\theta^2/\alpha}\,uQ_\alpha\}.
 \end{aligned}
 \tag{HO23}
\]
Every other coefficient here is infinitesimal by (HO2),(HO21),
because $\theta\sqrt{m/\alpha}=O(\sqrt x)$.
The complete bounded Gram on (HO19) therefore has strictly positive
diagonal residue
\[
             \operatorname{diag}(1,1,1,1,1,1/2,1/2).
 \tag{HO24}
\]
This proves the reciprocal chart for the full parent and every state.

\paragraph{All-tests join and receipt.}
The inherited small-$L$, bounded-$\eta$ remote theorem covers the
complement of both-infinitesimal testing angles on this original
physical source domain. Its exact source openings and literal beta
zero remain available. The original exact sign gauges, literal axes,
homogeneous state extension and real/imaginary decomposition join
as in IC1, RO and JO. All roots, radicals and inverses above are
pure-field or separated; every bound is a product of a source bound
and a testing bound.

The separate checker records 28 exact assertions, including the
full-parent scalar substitution, all three complete first normals,
the two full-event maps, the harmonic completion, the exact reference
kernel cancellation, the whole reciprocal remainder, its two strict
soft pivots and the retained-unit denominator factorizations.
Those checks support the finite bounds and circuit allocation proved
above; they are not being used as a substitute for them.
This theorem removes precisely the simultaneous positive
infinitesimal $r,j$ opening with unit $\eta,\zeta$.

\begin{corollary}[The entire unit-$\eta,\zeta$ small-$L$ region]
Within the original physical native small-$L$ C1.4 source, suppose
$x>0$ is infinitesimal and $\eta,\zeta$ are positive bounded units.
Then the full parent is tensor SOS at every independent projective
test and every state, with all original source intersections retained.
\end{corollary}
\begin{proof}
The positive actual gap $j=\eta(1+xh)-1+r$ is bounded.
If $j$ is a unit, then $j_0=j-\eta xh$ is a positive unit,
so RO applies, including either $r$ opening. If $j$ is infinitesimal,
then $1-r\sim\eta$ is a unit. A unit $r$ is exactly the JO domain;
an infinitesimal $r$ is exactly the HO domain. These alternatives
exhaust the stated source, require no comparison of $j$ with $x$,
and each supplies its complete local, remote and endpoint join.
\end{proof}
